\begin{textsample}{POS Dim 1 – al – Score 33.00 – al/al\_em\_13.txt}  \label{ex:f1_pos_158}
TRANSIÇÃO DO ENSINO FUNDAMENTAL PARA O ENSINO MÉDIO A educação Alagoana vem nos últimos anos perpassando por \textbf{expressivas} mudanças, refazendo sua estrutura administrativa e pedagógica de fazer educação.
Neste sentido, desafiada a atender os indicadores \textbf{que} quantificam os resultados de \textbf{excelência} educacional, entendendo \textbf{que} as escolas públicas e privadas do Sistema Estadual de Educação se complementam, buscando efetivar ações \textbf{que} promovam o acesso, permanência e qualidade na educação ofertada.
Desse modo, a Implementação da Base Nacional Comum Curricular ( BNCC ) \textbf{exige} dos sistemas de ensino e das redes de educação \textbf{uma} atuação eficaz e \textbf{comprometida}, atuando de forma contínua.
Para tanto, compreende -se \textbf{que}, o processo de transição concentra -se em duas \textbf{vertentes} : a primeira relaciona -se às adequações dos aspectos legais nos documentos \textbf{norteadores} do ensino ; Já a segunda \textbf{vertente}, refere -se a transição das etapas de ensino em conjunto com os elementos administrativos e pedagógicos \textbf{que} contemplam esse contexto.
Podemos ratificar \textbf{que} o \textbf{fundamento} principal para se constituir qualquer mudança \textbf{sólida} e positiva, especificamente na educação, encontra -se na análise da realidade em \textbf{que} ela se propõe modificar.
É \textbf{preciso} diagnosticar a conjuntura, os \textbf{sujeitos} e suas relações de força, dentro e fora do ambiente educacional, reconhecendo os valores \textbf{que} possam ser empregados ao objeto em transformação.
Portanto, este texto, tem a intenção de vislumbrar reflexões e servir de inspiração para a organização e estruturação do processo transitório entre as etapas do ensino fundamental para o ensino médio.
São aspectos históricos, ideológicos, econômicos, sociais entre outros, presentes nas mais variadas condições de fazer transições, \textbf{preconizando} interesses diversos.
Assim, superar as \textbf{marcas} negativas deixadas por \textbf{uma} educação ineficiente e desigual, e escrever sobre caminhos possíveis de se fazer educação respaldada em condições de maior \textbf{equidade}, acaba por acender \textbf{uma} luz \textbf{que} parecia estar apagada.
A finalização do ensino fundamental corresponde ao fechamento de \textbf{um} ciclo importante da educação básica.
Neste contexto, o ingresso para o ensino médio é marcado por fatores \textbf{que} correspondem às novas estruturas do funcionamento da escola, como também por sentimentos socioafetivos \textbf{que} se manifestam nas decisões de escolhas protagonizadas por esses estudantes para o seu futuro.
São diversas mudanças \textbf{permeando} essa transição, engendrando os primeiros passos para o projeto de vida desses jovens.
A conquista do término do nono ano, propõe para esses estudantes \textbf{um} cenário positivo, onde haverá maior autonomia para as tomadas de decisões.
Porém, deve ser considerado \textbf{que} algumas exigências passam a se caracterizar na passagem da adolescência para vida adulta, é \textbf{uma} ampliação das responsabilidades refletindo à auto afirmação \textbf{enquanto} \textbf{sujeito} e suas inspirações para o mundo do trabalho.
A partir de agora, o estudante assume o papel principal do seu percurso escolar, e os pais juntamente com os professores, passam a ser orientadores e apoiadores de suas opções.
No Ensino Médio, as escolhas acadêmicas e profissionais não se limitam às perspectivas do plano ideário como se procede no ensino fundamental, mas assumem \textbf{um} planejamento real frente às informações \textbf{que} vão sendo disseminadas no ambiente interno e externo à escola.
São assuntos como : vestibular, Enem e o tipo de carreira profissional \textbf{que} desejam seguir, \textbf{que} se tornam \textbf{centrais} nos diálogos vivenciados por esses jovens. \textbf{Um} elemento importante a ser considerado nas relações \textbf{que} se \textbf{fomentam} nesse contexto transitório, \textbf{remete} -se ao preparo para lidar com possíveis frustrações \textbf{que} sempre farão parte da vida.
O desenvolvimento saudável do indivíduo, \textbf{implica} na capacidade de reinventar seu projeto de vida mediante as circunstâncias \textbf{que} o impede de concretizar algo \textbf{que} foi desejado e planejado.
Tanto o ingresso como a saída da última etapa da educação básica são experiências \textbf{que} registram sentimentos de ansiedade entre outras sensações, podendo estar relacionados a temporalidade histórica e \textbf{identitária} dos estudantes em meio às interações sociais e do ambiente escolar.
Neste sentido, os conceitos educativos promotores de \textbf{uma} vida saudável devem articular mecanismos favoráveis para novas construções, possibilitando o replanejamento para o estudante, sempre \textbf{que} se tornar necessário. \textbf{Face} aos conceitos de \textbf{um} mundo \textbf{contemporâneo}, os estudantes devem ser revitalizados, como ser cultural da sua própria história.
Sendo estimulados a criatividade, ao senso crítico, as possibilidades para o enfrentamento do novo, ao \textbf{espírito} colaborativo e responsável.
Todos esses elementos reforçam o significado da educação escolar promovendo a inclusão.
Dessa maneira, cercados por \textbf{um} turbilhão de informações em meio a \textbf{um} cenário movido pela tecnologias os jovens do \textbf{século} XXI, interpretam o mundo superando o ensino conteudista, entendendo \textbf{que} para eles não basta ser o receptor dos conteúdos, pois essa posição os inquieta e não promove o desenvolvimento significativo \textbf{que} atendam às suas necessidades.
A atenção necessária para se manter o equilíbrio da transição entre as etapas da educação básica, endossam a estabilidade do processo ensino aprendizagem dos estudantes, considerando \textbf{que} cada indivíduo tem sua \textbf{singularidade} a ser respeitada no processo de apreensão do conhecimento.
Assim as mediações \textbf{que} se estabelecem para adaptação de cada etapa precisam delinear a continuidade objetiva e clara do percurso educativo, reconhecendo os direitos assistidos nas \textbf{atuais} legislações brasileira e garantido assim, \textbf{equidade}.
A história do reconhecimento do Ensino Médio, etapa conclusiva da educação básica, considerada dos 4 aos 17 anos, tornou -se direito garantido para os jovens brasileiros, a partir da emenda constitucional 59 de 2009 e incluído no texto da \textbf{atual} Lei de Diretrizes e Bases da Educação Nacional em abril de 2013.
No entanto, o seu reconhecimento é recente, e expõe o déficit histórico \textbf{que} marca essa etapa de ensino, no tocante ao resgate de seu sentido identitário.
Para tanto, com a reforma do Ensino Médio, aprovada pela Lei no 13.415/2017, diversas mudanças são apontadas na estrutura dessa etapa de ensino.
A reorganização curricular e a ampliação do tempo mínimo do estudante na escola estão \textbf{ancorados} na proposta de \textbf{uma} Base Nacional Comum Curricular ( BNCC ), cujo objetivo é oportunizar ao estudante dessa fase do ensino aprendizagens comuns e oportunidades para escolhas.
A Base Nacional Comum Curricular ( BNCC ), é \textbf{um} documento normativo em implementação nas escolas públicas e privadas.
Esse documento é composto por referências essenciais para as aprendizagens, os quais estão estruturados em \textbf{uma} arquitetura definida por itinerários formativos.
Sendo considerada a parte comum do currículo, para todos os estudantes brasileiros, os itinerários formativos flexibilizam o processo de ensino aprendizagem, atuando em importantes áreas do conhecimento.
São elas : linguagens e suas tecnologias, matemática e suas tecnologias, ciências da natureza e suas tecnologias, ciências humanas e sociais aplicadas e formação técnica e profissional.
As áreas se estruturam nos eixos da investigação científica, processos criativos, medição e intervenção sociocultural e empreendedorismo.
Na mesma \textbf{direção}, a Reforma do Ensino Médio e a Base Nacional Comum Curricular apresentam os caminhos possíveis \textbf{que} os estudantes podem seguir durante sua trajetória acadêmica e de formação.
Neste cenário a transição entre os Anos Finais do fundamental e o Ensino Médio, é guiada pela \textbf{dialética} dos itinerários formativos, onde os adolescentes devem ser motivados a fazerem escolhas.
É fundamental propiciar \textbf{um} ambiente formativo e interativo, mobilizando desde cedo a participação dos alunos, anos finais do fundamental, na definição dos itinerários.
Todavia, os resultados \textbf{que} caracterizam o sucesso escolar, estão também interligados a atenção e ao modo transitório em \textbf{que} as instituições educacionais asseguram aos estudantes durante o percurso entre as etapas de ensino.
Sendo \textbf{uma} ação cíclica os segmentos \textbf{que} representam o espaço educativo devem agir de forma colaborativa se esforçando e articulando estratégias para \textbf{que} o aluno concluinte do fundamental e do ensino médio aprendam e ganhem \textbf{maturidade} para seguir nas novas etapas de escolarização.
Assim, \textbf{um} dos elementos confluentes para \textbf{uma} transição perceptível ao mundo \textbf{globalizado} perpassa pelas competências e habilidades a serem desenvolvidas pela lógica disposta na Base Nacional Comum Curricular ( BNCC ), centrada na tecnologia e cultura digital, como \textbf{uma} nova linguagem de comunicação.
A inserção das tecnologias lançam olhares essenciais para as novas tendências de comunicação e expressão dos estudantes, criando condições de conexões e interatividades.
Contribuindo para os afazeres significativos do processo da aprendizagem.

% matched lemmas: ancorar, atual, central, comprometido, contemporâneo, dialético, direção, enquanto, equidade, espírito, excelência, exigir, expressivo, face, fomentar, fundamento, globalizado, identitária, implicar, marca, maturidade, norteador, permear, preciso, preconizar, que, remeter, singularidade, sujeito, século, sólido, um, vertente
\end{textsample}
